%==============================================================
\section{Cổng lượng tử và mạch lượng tử}
\label{sec:gates}
%==============================================================

\subsection{Cổng lượng tử là ma trận unitary}
\label{sec:unitary}

\begin{definition}[Ma trận unitary]
Ma trận vuông phức $U$ gọi là \textbf{unitary} nếu
\begin{equation}
  U^\dagger U = U U^\dagger = I,
  \label{eq:unitary}
\end{equation}
trong đó $U^\dagger = \overline{U}^{\mathsf T}$ là \textbf{liên hợp Hermite} (chuyển vị liên hợp) của
$U$: chuyển vị ma trận rồi lấy liên hợp phức từng phần tử.
\end{definition}

Một cổng lượng tử tác dụng lên $n$ qubit là một ma trận unitary kích thước $2^n \times 2^n$, và áp cổng
lên trạng thái chính là phép nhân ma trận: $\ket{\psi'} = U\ket{\psi}$. Điều kiện~\eqref{eq:unitary}
nói rằng \textbf{nghịch đảo của $U$ chính là $U^\dagger$}, nên luôn tính được trực tiếp. Ví dụ với cổng
Hadamard và cổng phase $S$:
\[
  H^\dagger H = \tfrac12\begin{pmatrix}1&1\\1&-1\end{pmatrix}\begin{pmatrix}1&1\\1&-1\end{pmatrix}
  = \tfrac12\begin{pmatrix}2&0\\0&2\end{pmatrix} = I, \qquad
  S^\dagger S = \begin{pmatrix}1&0\\0&-i\end{pmatrix}\begin{pmatrix}1&0\\0&i\end{pmatrix} = I.
\]
Cổng $S^\dagger$ chính là \code{sdg} trong Qiskit -- tên gọi lấy từ ``S-dagger''.

\paragraph{Vì sao cổng lượng tử bắt buộc phải unitary.} Một cổng phải biến trạng thái hợp lệ thành
trạng thái hợp lệ, tức tổng xác suất luôn bằng 1:
\begin{equation}
  \big\|U\ket{\psi}\big\|^2 = \bra{\psi}U^\dagger U\ket{\psi} = \braket{\psi|\psi} = 1.
\end{equation}
Đẳng thức này đúng với \emph{mọi} $\ket{\psi}$ khi và chỉ khi $U^\dagger U = I$. Vậy tính unitary không
phải một giả thiết tuỳ ý: unitary chính xác là tập các phép biến đổi tuyến tính \textbf{bảo toàn xác
suất}. Về mặt vật lý, một hệ lượng tử cô lập tiến hoá theo phương trình Schrödinger với nghiệm
$U = e^{-i\mathcal{H}t}$, trong đó Hamiltonian $\mathcal{H}$ là một ma trận Hermite; ma trận dạng này
luôn unitary. Thực hiện một cổng trên phần cứng chính là điều khiển Hamiltonian của hệ (bằng xung vi
sóng, laser, \dots) trong một khoảng thời gian xác định.

\begin{table}[h]
  \centering
  \caption{Các tính chất của ma trận unitary và ý nghĩa đối với mạch lượng tử.}
  \label{tab:unitary}
  \small
  \begin{tabularx}{\linewidth}{@{}lX@{}}
    \toprule
    \textbf{Tính chất} & \textbf{Ý nghĩa} \\
    \midrule
    Bảo toàn tích trong: $\braket{U\phi|U\psi} = \braket{\phi|\psi}$ & giữ nguyên ``góc'' giữa các
      trạng thái; hai trạng thái trực giao vẫn trực giao sau cổng \\
    Khả nghịch: $U^{-1} = U^\dagger$ & không cổng nào ``xoá'' thông tin; \code{qc.inverse()} luôn tồn tại \\
    Các cột (và hàng) trực chuẩn & cột thứ $j$ là $U\ket{j}$; ảnh của một cơ sở vẫn là một cơ sở \\
    Trị riêng có môđun 1: $\lambda = e^{i\theta}$ & cổng chỉ ``quay pha'' chứ không co giãn vector \\
    Đóng với phép nhân và tích tensor & ghép cổng nối tiếp ($UV$) hay song song ($U\otimes V$) vẫn được
      unitary, nên cả mạch (không có phép đo) là một unitary \\
    Viết được dưới dạng $U = e^{iA}$, $A$ Hermite & nguồn gốc của các cổng quay
      $R_P(\theta) = e^{-i\theta P/2}$ ở Mục~\ref{sec:1q} \\
    \bottomrule
  \end{tabularx}
\end{table}

\paragraph{So sánh với logic cổ điển.} Cổng AND cổ điển \textbf{không} khả nghịch: đầu ra 0 có thể đến
từ \code{00}, \code{01} hoặc \code{10}, nên thông tin về đầu vào bị mất. Vì vậy AND không thể là một
cổng lượng tử trực tiếp. Muốn tính một hàm cổ điển $f$ trong mạch lượng tử, phải viết lại nó ở dạng
khả nghịch -- cổng Toffoli (một bản AND khả nghịch, Mục~\ref{sec:mq}) hoặc bit-flip oracle
$U_f\ket{x}\ket{y} = \ket{x}\ket{y\oplus f(x)}$ ở Mục~\ref{sec:algorithms}.

\paragraph{Không phải lệnh nào trong mạch cũng unitary.} Phép đo (\code{measure}) và đặt lại qubit
(\code{reset}) làm mất thông tin nên không khả nghịch, và do đó không phải cổng unitary. Đây là lý do
\code{Operator(qc)} và \code{Statevector(qc)} chỉ dùng được với mạch không chứa hai lệnh này. Qiskit
cũng cho phép tạo cổng trực tiếp từ một ma trận bằng \code{UnitaryGate(M)}, và kiểm tra một ma trận
bằng \code{Operator(M).is\_unitary()}.

\paragraph{Số tham số và bài toán synthesis.} Một unitary trên $n$ qubit có $4^n$ tham số thực. Với
một qubit là 4 tham số; bỏ global phase còn 3 -- khớp với ba góc của cổng tổng quát
$U(\theta,\varphi,\lambda)$ ở phương trình~\eqref{eq:u3}. Với hai qubit là 16 tham số, và định lý KAK
cho biết mọi unitary hai qubit cần tối đa 3 cổng CX. Số tham số tăng theo hàm mũ khiến việc phân tích
một unitary tuỳ ý thành chuỗi cổng cơ bản -- bài toán \emph{synthesis} -- trở nên rất khó khi số qubit
tăng (Mục~\ref{sec:synthesis}).

\subsection{Cổng một qubit}
\label{sec:1q}

\paragraph{Nhóm Pauli.}
\begin{equation}
  X = \begin{pmatrix}0&1\\1&0\end{pmatrix}, \qquad
  Y = \begin{pmatrix}0&-i\\i&0\end{pmatrix}, \qquad
  Z = \begin{pmatrix}1&0\\0&-1\end{pmatrix}.
\end{equation}
$X$ là ``NOT lượng tử'' (bit flip, $\ket{0}\leftrightarrow\ket{1}$); $Z$ là phase flip
($\ket{1}\to-\ket{1}$); $Y = iXZ$ làm cả hai. Trên Bloch sphere chúng là các phép quay $180^\circ$
quanh trục tương ứng, và $X^2 = Y^2 = Z^2 = I$.

\paragraph{Hadamard.}
\begin{equation}
  H = \frac{1}{\sqrt2}\begin{pmatrix}1&1\\1&-1\end{pmatrix}, \qquad
  H\ket{0} = \ket{+}, \quad H\ket{1} = \ket{-}, \quad H^2 = I.
\end{equation}
$H$ chuyển qua lại giữa cơ sở $Z$ và cơ sở $X$, nên ``thêm $H$ trước khi đo'' chính là đo trong cơ
sở $X$. Đẳng thức $HZH = X$ (và $HXH = Z$) được dùng liên tục trong thiết kế và tối ưu mạch.

\paragraph{Cổng phase.}
\begin{equation}
  S = \begin{pmatrix}1&0\\0&i\end{pmatrix}, \qquad
  T = \begin{pmatrix}1&0\\0&e^{i\pi/4}\end{pmatrix}, \qquad
  P(\lambda) = \begin{pmatrix}1&0\\0&e^{i\lambda}\end{pmatrix},
  \qquad S^2 = Z,\ T^2 = S.
\end{equation}
Các cổng này chỉ thêm phase cho thành phần $\ket{1}$, tức quay quanh trục $Z$. Vì không đổi
$|\alpha|, |\beta|$ nên đo ngay sau cổng phase không thấy khác biệt gì; phase chỉ lộ ra khi kết hợp
với $H$. Cổng $T$ đặc biệt quan trọng ở Mục~\ref{sec:synthesis}: nó là cổng ``đắt'' nhất trên máy
tính lượng tử chịu lỗi.

\paragraph{Cổng quay.} Với góc tuỳ ý, dùng các cổng rotation
$R_P(\theta) = e^{-i\theta P/2} = \cos\tfrac{\theta}{2}I - i\sin\tfrac{\theta}{2}P$ với
$P \in \{X, Y, Z\}$, ví dụ
\begin{equation}
  R_Y(\theta) = \begin{pmatrix}\cos\frac\theta2 & -\sin\frac\theta2\\[2pt]
                               \sin\frac\theta2 & \cos\frac\theta2\end{pmatrix}, \qquad
  R_Z(\theta) = \begin{pmatrix}e^{-i\theta/2}&0\\0&e^{i\theta/2}\end{pmatrix}.
\end{equation}
Chú ý hệ số $\theta/2$: phải quay $\theta = 2\pi$ mới quay về đúng trạng thái ban đầu kể cả phase.
Ví dụ $R_X(\pi) = -iX$ khác $X$ đúng một global phase.

\paragraph{Cổng tổng quát.} Mọi cổng một qubit đều là trường hợp riêng (sai khác global phase) của
\begin{equation}
  U(\theta,\varphi,\lambda) = \begin{pmatrix}
    \cos\frac\theta2 & -e^{i\lambda}\sin\frac\theta2\\[2pt]
    e^{i\varphi}\sin\frac\theta2 & e^{i(\varphi+\lambda)}\cos\frac\theta2\end{pmatrix}.
  \label{eq:u3}
\end{equation}
Ba tham số thực khớp với ba bậc tự do của $SU(2)$ -- kết quả này là nền tảng của synthesis một qubit
ở Mục~\ref{sec:synthesis}.

\begin{table}[h]
  \centering
  \caption{Tóm tắt các cổng một qubit. Ký hiệu $\mathrm{diag}(a,b)$ là ma trận đường chéo.}
  \label{tab:1q}
  \small
  \begin{tabular}{@{}llll@{}}
    \toprule
    \textbf{Cổng} & \textbf{Qiskit} & \textbf{Quay trên Bloch sphere} & \textbf{Ghi nhớ} \\
    \midrule
    $X$ & \code{qc.x(0)} & $180^\circ$ quanh $X$ & NOT lượng tử \\
    $Y$ & \code{qc.y(0)} & $180^\circ$ quanh $Y$ & bit flip + phase flip \\
    $Z$ & \code{qc.z(0)} & $180^\circ$ quanh $Z$ & phase flip \\
    $H$ & \code{qc.h(0)} & $180^\circ$ quanh $(X+Z)/\sqrt2$ & tạo superposition \\
    $S = \mathrm{diag}(1, i)$ & \code{qc.s(0)} & $90^\circ$ quanh $Z$ & $S^2 = Z$ \\
    $T = \mathrm{diag}(1, e^{i\pi/4})$ & \code{qc.t(0)} & $45^\circ$ quanh $Z$ & $T^2 = S$ \\
    $P(\lambda)$ & \code{qc.p(lam, 0)} & $\lambda$ quanh $Z$ & tổng quát của $Z, S, T$ \\
    $R_X, R_Y, R_Z$ & \code{qc.rx/ry/rz(th, 0)} & $\theta$ quanh trục tương ứng & $R_Y$: $P(1)=\sin^2\frac\theta2$ \\
    \bottomrule
  \end{tabular}
\end{table}

\paragraph{Thứ tự cổng và thứ tự nhân ma trận.} Mạch được vẽ theo thời gian từ trái sang phải, nhưng
ma trận tương ứng được nhân từ phải sang trái. Mạch ``$H$ rồi $X$'' có ma trận $U = X \cdot H$, vì
$X(H\ket{\psi}) = (XH)\ket{\psi}$. Nghịch đảo của cả mạch (\code{qc.inverse()}) đảo ngược thứ tự các
cổng và lấy liên hợp Hermite từng cổng ($T^\dagger$ là \code{tdg}, $S^\dagger$ là \code{sdg}).

\subsection{Cổng nhiều qubit}
\label{sec:mq}

\paragraph{CNOT (CX).} Cổng hai qubit quan trọng nhất: có một qubit \emph{control} và một qubit
\emph{target}; nếu control bằng 1 thì lật target:
\begin{equation}
  \mathrm{CX}\,\ket{c}\ket{t} = \ket{c}\ket{t \oplus c}.
\end{equation}
CX là cổng tạo \emph{tương quan} giữa hai qubit, và cũng là cổng khó thực hiện nhất trên phần cứng.
Vì vậy \textbf{số cổng CX là thước đo chi phí chính} của một mạch -- ý này sẽ xuất hiện lại ở mọi phần
của Mục~\ref{sec:compilation}.

\paragraph{Các cổng khác.} Bảng~\ref{tab:mq} liệt kê các cổng nhiều qubit thường dùng. Một quan hệ
đáng nhớ: kẹp $H$ lên target trước và sau CX thì được CZ, tức $\mathrm{CZ} = (I\otimes H)\,\mathrm{CX}\,(I\otimes H)$
(viết theo thứ tự qubit control $\otimes$ target). Một cổng bất kỳ có thể được biến thành cổng
controlled bằng \code{.to\_gate().control(n)}.

\begin{table}[h]
  \centering
  \caption{Các cổng nhiều qubit thường dùng.}
  \label{tab:mq}
  \small
  \begin{tabular}{@{}lll@{}}
    \toprule
    \textbf{Cổng} & \textbf{Qiskit} & \textbf{Ý nghĩa} \\
    \midrule
    CX / CNOT & \code{qc.cx(c, t)} & control $=1$ thì lật target \\
    CZ & \code{qc.cz(a, b)} & cả hai bằng 1 thì thêm dấu $-$; \textbf{đối xứng} \\
    SWAP & \code{qc.swap(a, b)} & đổi chỗ trạng thái hai qubit; $=3$ CX \\
    Toffoli (CCX) & \code{qc.ccx(c1, c2, t)} & hai control; đủ để làm mọi hàm logic cổ điển khả nghịch \\
    CH, CRX, CP, \dots & \code{qc.ch}, \code{qc.crx}, \code{qc.cp} & bản controlled của cổng một qubit \\
    \bottomrule
  \end{tabular}
\end{table}

\subsection{Entanglement}

\paragraph{Trạng thái Bell.} Chỉ cần hai cổng để tạo ra entanglement:
\begin{equation}
  \ket{00} \xrightarrow{\;H \text{ trên } q_0\;} \tfrac{1}{\sqrt2}\big(\ket{00}+\ket{01}\big)
  \xrightarrow{\;\mathrm{CX}_{0\to1}\;} \tfrac{1}{\sqrt2}\big(\ket{00}+\ket{11}\big) = \ket{\Phi^+}.
\end{equation}

\begin{figure}[h]
  \centering
  \begin{quantikz}
    \lstick{$q_0: \ket{0}$} & \gate{H} & \ctrl{1} & \meter{} \\
    \lstick{$q_1: \ket{0}$} & & \targ{} & \meter{}
  \end{quantikz}
  \caption{Mạch tạo trạng thái Bell $\ket{\Phi^+}$. Kết quả đo chỉ có \code{00} hoặc \code{11},
  mỗi chuỗi xác suất $1/2$.}
  \label{fig:bell}
\end{figure}

Kết quả đo chỉ có \code{00} và \code{11}, không bao giờ có \code{01} hay \code{10}: hai qubit luôn
cho cùng kết quả. Bốn trạng thái Bell
\begin{equation}
  \ket{\Phi^\pm} = \tfrac{1}{\sqrt2}\big(\ket{00}\pm\ket{11}\big), \qquad
  \ket{\Psi^\pm} = \tfrac{1}{\sqrt2}\big(\ket{01}\pm\ket{10}\big)
\end{equation}
trực giao nhau và tạo thành \textbf{Bell basis} của không gian hai qubit. Mở rộng ra $n$ qubit, ta có
trạng thái GHZ $\tfrac{1}{\sqrt2}(\ket{0\dots0}+\ket{1\dots1})$, tạo bằng một cổng $H$ và một chuỗi
CX.

\begin{definition}[Separable và entangled]
Trạng thái hai qubit là \textbf{separable} nếu viết được thành $\ket{\psi_A}\otimes\ket{\psi_B}$;
ngược lại nó là \textbf{entangled}.
\end{definition}

\paragraph{Đo lường entanglement.} Lấy \emph{partial trace} để bỏ qubit $B$, ta được ma trận mật độ
(density matrix) $\rho_A = \Tr_B\ket{\psi}\bra{\psi}$ của qubit $A$. Khi đó:
\begin{itemize}[noitemsep]
  \item \textbf{Purity} $\Tr(\rho_A^2)$ bằng 1 nếu $A$ ở pure state, bằng $0.5$ nếu entangled tối đa.
  \item \textbf{Entropy von Neumann} $S(\rho_A) = -\Tr(\rho_A\log_2\rho_A)$ bằng 0 nếu không
    entangled, bằng 1 nếu entangled tối đa.
\end{itemize}
Với $\ket{\Phi^+}$, giá trị kỳ vọng tính bằng \code{EstimatorV2} cho $\langle ZZ\rangle = +1$ và
$\langle XX\rangle = +1$ (tương quan hoàn hảo trong \emph{cả hai} cơ sở -- điều mà tương quan cổ điển
không làm được), trong khi $\langle ZI\rangle = 0$: riêng từng qubit thì hoàn toàn ngẫu nhiên. Toàn
bộ thông tin nằm ở mối tương quan giữa hai qubit.

\subsection{Tính universal của tập cổng}

Một tập cổng gọi là \textbf{universal} nếu mọi unitary đều biểu diễn được (chính xác hoặc xấp xỉ tuỳ
ý) bằng các cổng trong tập. Hai kết quả kinh điển \cite{nielsen2010, barenco1995}:
\begin{itemize}[noitemsep]
  \item Các cổng một qubit cùng với CX là universal \emph{chính xác}. Đây là cơ sở của các basis
    gates trên phần cứng NISQ, ví dụ $\{\code{rz}, \code{sx}, \code{x}, \code{cx}\}$.
  \item Tập rời rạc Clifford+T $= \{H, S, \mathrm{CX}, T\}$ là universal \emph{xấp xỉ}: mọi unitary
    xấp xỉ được tới sai số $\epsilon$ tuỳ ý. Đây là tập cổng của máy tính lượng tử chịu lỗi.
\end{itemize}
Sự khác biệt giữa hai loại universal này dẫn tới hai bài toán synthesis rất khác nhau ở
Mục~\ref{sec:synthesis}.

\subsection{Quy trình làm việc với Qiskit 2.x}
\label{sec:workflow}

Mọi bài toán trong báo cáo đều đi qua năm bước: (1) \textbf{xây mạch} bằng \code{QuantumCircuit};
(2) \textbf{xem mạch} bằng \code{draw()}, \code{depth()}, \code{count\_ops()}; (3) \textbf{đo} bằng
\code{measure()} hoặc \code{measure\_all()}; (4) \textbf{chạy} bằng một trong ba cách ở
Bảng~\ref{tab:primitives}; (5) \textbf{transpile} về tập cổng và coupling map của thiết bị.

\begin{lstlisting}[caption={Quy trình tối thiểu: tạo trạng thái Bell, transpile và lấy mẫu.}]
from qiskit import QuantumCircuit, transpile
from qiskit_aer import AerSimulator
from qiskit_aer.primitives import SamplerV2

qc = QuantumCircuit(2)
qc.h(0); qc.cx(0, 1)
qc.measure_all()                      # creates a register named "meas"

backend = AerSimulator()
tqc = transpile(qc, backend, optimization_level=2)
res = SamplerV2().run([tqc], shots=1000).result()
print(res[0].data.meas.get_counts())  # ~ {'00': 500, '11': 500}
\end{lstlisting}

\begin{table}[h]
  \centering
  \caption{Ba cách lấy kết quả từ mạch trong Qiskit 2.x.}
  \label{tab:primitives}
  \small
  \begin{tabular}{@{}llll@{}}
    \toprule
    \textbf{Cách} & \textbf{Cho ra} & \textbf{Mạch có \code{measure}?} & \textbf{Dùng khi} \\
    \midrule
    \code{Statevector(qc)} & amplitude chính xác & \textbf{không} được có & học, debug mạch nhỏ \\
    \code{SamplerV2} & counts từ các shot & \textbf{bắt buộc} có & mô phỏng kết quả thật \\
    \code{EstimatorV2} & giá trị kỳ vọng $\langle O\rangle$ & không cần & VQE, đo tương quan \\
    \bottomrule
  \end{tabular}
\end{table}

\begin{warnbox}[Các lỗi hay gặp khi mới dùng Qiskit 2.x]
(1) Đọc nhầm thanh ghi: \code{measure\_all()} tạo thanh ghi \code{meas}, không phải \code{c}.
(2) Gọi \code{Statevector(qc)} khi mạch đã có lệnh đo.
(3) Quên \code{transpile} trước \code{backend.run()}.
(4) Nhầm thứ tự bit: qubit 0 là ký tự bên phải.
(5) So sánh cổng bằng \code{==} trong khi chỉ cần \code{equiv} (khác nhau ở global phase).
Ngoài ra, nhiều tài liệu trên mạng dùng API Qiskit 0.x đã bị gỡ bỏ (\code{execute},
\code{Aer.get\_backend}); API tương ứng hiện nay là các primitive \code{SamplerV2}/\code{EstimatorV2}
và \code{AerSimulator}.
\end{warnbox}
